\section{Adjustment as a Substitute for Enforcement}
\label{sec:enforcementfrontier}
The comparison of compulsory terms is not the only way in which adjustment and enforcement interact. An agent who values continuation more highly requires less protection against a valuable surrender option. This channel can be stated in the full model, with endogenous consumption, financial positions, repeated preference adjustment, and stopping. It does not require an irreversible-upgrade subclass or exogenous consumption.

\subsection{Initial-state and uniform enforcement requirements}
Fix a mandate and compulsory term. Absorb the operating benefit into a policy's pre-fee payoff $B(p)$ and let $H(p)$ be discounted elective surrender. Let $W_\infty$ be the best payoff among policies that never surrender electively. Forced discharge remains part of this problem. Assume that the class with $H=0$ is nonempty and that all values are finite. Define
\begin{equation}
 F^{\rm init}=\max\left\{0,\sup_{p:H(p)>0}
                   \frac{B(p)-W_\infty}{H(p)}\right\}.
 \label{eq:initialthreshold}
\end{equation}
The supremum over an empty set is omitted. A threshold above the legal fee ceiling is a property of the economic environment, not a reason to change the ceiling silently.

\begin{theorem}[Enforcement frontiers in the full dynamic model]
\label{thm:enforcementfrontier}
For a fixed initial condition, $W(F)=W_\infty$ if and only if $F\geq F^{\rm init}$. For $F>F^{\rm init}$, every attained exact response has $H=0$. At the threshold, an exact response with positive surrender can tie a no-surrender policy.

For implementation from every live state at every date when surrender is permitted, let $V^r_{\infty,n}$ be the no-elective-surrender value and $G_n$ the surrender settlement before the fee. The uniform value-preservation threshold, and the infimum of fees eliminating every elective-stop tie, is
\begin{equation}
 F^{\rm unif}_r=\sup_{n\geq m,\ s\ {\rm live}}
                    [G_n(s)-V^r_{\infty,n}(s)]_+.
 \label{eq:uniformthreshold}
\end{equation}
Suppose adjustment enlarges the feasible operating controls without changing the settlement payoff, discounting, fee incidence in private utility, or forced-discharge technology. Then
\begin{equation}
 V^A_{\infty,n}(s)\geq V^0_{\infty,n}(s),\qquad
 F^{\rm unif}_A\leq F^{\rm unif}_0.
 \label{eq:uniformordering}
\end{equation}
These conclusions hold for the finite dynamic model and, under the usual dynamic-programming and comparison hypotheses, for a stopped controlled diffusion. The continuous-time version replaces dates by permitted stopping times and requires an obstacle-compatible value function. No restriction to exogenous consumption or a single preference change is used.
\end{theorem}

Uniform and initial-state enforcement are different economic requirements. A single reachable state can matter for exact implementation even if its probability is extremely small. Conversely, a state outside the support of the initial operating process cannot determine an initial-state certificate merely because it appears on a rectangular grid. Full-support conditions and an attained no-surrender optimum can make the uniform obstacle bound necessary for initial implementation as well. Without those additional conditions, \eqref{eq:uniformordering} does not by itself rank $F^{\rm init}$ across financial mandates or adjustment regimes. We therefore calculate the initial-state thresholds separately rather than infer them from policy-set inclusion.

\subsection{A verified comparison with endogenous stopping}
For the positive initial mandate and term one, the no-surrender value is obtained by disabling elective surrender, not by imposing an arbitrarily large numerical fee. A lower threshold endpoint is certified by a strictly positive value advantage of surrender over this no-surrender problem. At the upper endpoint, forward reachability through a superset of all Bellman-optimal actions finds no elective-stop action. Thus the zero-surrender conclusion is structural: a small computed probability is not rounded to zero.

\input{revisions/2026-09-19-r15-referee-response/paper/table_enforcement}

The same arrays also permit calculation of the uniform obstacle requirement over all live states and permitted dates. The uniform fee need not lie in the baseline administrative interval $[0,1]$. The uniform value-preservation thresholds are about $1.991$ and $2.512$, respectively; both exceed the administrative ceiling one. They are diagnostics of state-wide implementation, not additional contracts inserted into the procurement comparison. Their difference from the initial-state requirements is informative about exposure to unfavorable operating states. At the common fee $0.765$, the adjusted initial problem has no elective surrender under any exact response, while the no-adjustment problem has a strictly valuable surrender option. The main comparison is therefore a substitution between private continuation value and enforcement, not a purported invariant difference in compulsory terms.

Under binding participation and $H=0$, purchaser surplus is
\begin{equation}
 \Pi=bA+W_\infty-G_0-C(F/\zeta)-D(m).
 \label{eq:nosurrenderprofit}
\end{equation}
If two no-surrender contracts have the same term and duration, the regime with the larger continuation value and smaller implementation fee is weakly preferred for every common nondecreasing capacity cost, regardless of who initially pays that cost or receives the surrender fee. Strict improvement of either component gives strict preference when the corresponding cost or payoff comparison is strict. The calculation reports the component comparisons separately; it does not assume exact equality of rounded duration numbers.

More generally, on a fixed response face with moments $(B,A,H)$, binding participation gives
\[
 \Pi=bA+B-G_0-(1-\rho_F)FH-C(F/\zeta)-D(m).
\]
For $0\leq\rho_F\leq1$, a larger fee weakly reduces surplus while the same response face remains selected. An efficient fee is therefore driven toward a response boundary, whereas a compulsory term changes the set of response boundaries. This explains why a coarse fee menu can create a term-length difference that disappears when an inexpensive fee just beyond a surrender boundary becomes available. It also explains why fee incidence can change the selected face even though it does not change the agent's private decision problem.

\subsection{Continuous institutional counterfactuals}
\label{sec:continuousinstitutions}
The central counterfactuals use the same fee continuum and the same dynamic arrays as the baseline. We change fee receipts, capacity curvature, and enforcement efficiency. For $C(E)=cE^p$, $p\geq1$, a cell midpoint $v$ gives the globally valid tangent
\[
 C(F/\zeta)\geq c\zeta^{-p}
     \{pv^{p-1}F+(1-p)v^p\}.
\]
The agent's share of this tangent enters participation; the purchaser's share enters the upper objective. A negative tangent intercept is not a negative true cost and is allowed only as a conservative lower bound. The supplement gives the exact coefficient reconstruction and full-cover checks.

\input{revisions/2026-09-19-r15-referee-response/paper/table_continuous_institutions}

Cost incidence admits a stronger conclusion than another numerical row. The canonical value upper bound at fee zero is strictly below the outside value for every term and mandate. Since private value is nonincreasing in the surrender fee, participation binds throughout the entire baseline continuum, under either assignment of nonnegative capacity cost. Proposition~\ref{prop:capacity} therefore proves incidence neutrality on the continuum, not just at sampled prices. The proof does not cover payments that change managed wealth.

The normalized primitives make this a theory and computational benchmark, not an estimated institutional treatment effect. The general ordering in Theorem~\ref{thm:enforcementfrontier} is primitive and transferable; the magnitudes and the initial-state threshold comparison are quantitative statements about the explicitly recorded economy. Wealth-qualified output is contractible only under an institution that observes the relevant account value. Otherwise it is a diagnostic feature rather than a permissible payment condition. No private-type reporting mechanism is part of the purchaser's problem.
